% GENERATED FILE. Edit canonical lessons, then run npm run generate:books.
% canonical-source-hash: fnv1a64:aa4c9409d4d0fbee
% canonical-lessons: PA-C316-larna, PA-C316-jhagra, PA-C316-maarna, PA-C316-galh, PA-C316-sulha

\chapter{Trouble in the Lane}
\label{ch:pa-a2-trouble-in-the-lane}
\hlchaptermodality{316}

\begin{chapteropening}
\textbf{By the end of this chapter:} \emph{I can say who quarrelled and how it ended.}

Complete the last lesson of chapter 316: I can say who quarrelled and how it ended.
\end{chapteropening}

\section[\texorpdfstring{\pa{ਲੜਨਾ} (laṛnā)}{ਲੜਨਾ (laṛnā)}]{\pa{ਲੜਨਾ} (laṛnā) — to fight}

\label{lesson:PA-C316-larna}



\begin{quote}
Before the new one: say the Punjabi for a piece of jewellery, then the Punjabi for a life.
\end{quote}

\subsection*{You'll want to know: \pa{ਲੜਨਾ}}
\textbf{\pa{ਲੜਨਾ}} — \emph{laṛnā} — \textquotedblleft{}to fight\textquotedblright{}.

\begin{quote}
\pa{ਕੱਲ੍ਹ} \pa{ਗਲੀ} \pa{ਵਿੱਚ} \pa{ਦੋ} \pa{ਮੁੰਡੇ} \pa{ਲੜੇ।}
\end{quote}

\begin{quote}
\emph{kall galī vich do muṇḍe laṛe.}
\end{quote}

\begin{quote}
Yesterday two boys fought in the lane.
\end{quote}

Past (I …): \pa{ਮੈਂ} \pa{ਆਪਣੇ} \pa{ਦੋਸਤ} \pa{ਨਾਲ} \pa{ਲੜਿਆ}, \emph{maiṁ āpṇe dost nāl laṛiā}.

\begin{cousinweb}
\textbf{\pa{ਲੜਨਾ}} (\emph{laṛnā}) is inherited from a reconstructed \emph{*laḍ}, \textquotedblleft{}fight\textquotedblright{}; CDIAL lists Punjabi \emph{laṛṇā}, \textquotedblleft{}to fight\textquotedblright{}, under it (CDIAL 10920).
\end{cousinweb}

\begin{grammarlens}[title={What you've built}]
One more word for saying what happened.
\end{grammarlens}

\subsection*{Your turn}
\begin{itemize}
\raggedright
  \item \emph{Say it:} \emph{laṛnā}
  \item \emph{Say it:} \emph{laṛnā}, once more
  \item \emph{Say it:} \emph{kall galī vich do muṇḍe laṛe.} — Yesterday two boys fought in the lane.
  \item \emph{Recall:} say the Punjabi for a piece of jewellery, then the Punjabi for a life, then say \emph{laṛnā} again
\end{itemize}

\begin{tcolorbox}[breakable,colback=teal!4,colframe=teal!35!black,title={Before you move on}]
What does \pa{ਲੜਨਾ} mean? (\textquotedblleft{}To fight\textquotedblright{}.) Say it once more.
\end{tcolorbox}
\section[\texorpdfstring{\pa{ਝਗੜਾ} (jhagṛā)}{ਝਗੜਾ (jhagṛā)}]{\pa{ਝਗੜਾ} (jhagṛā) — a quarrel}

\label{lesson:PA-C316-jhagra}



\begin{quote}
Before the new one: say the Punjabi for a life, then the Punjabi for to fight.
\end{quote}

\subsection*{You'll want to know: \pa{ਝਗੜਾ}}
\textbf{\pa{ਝਗੜਾ}} — \emph{jhagṛā} — \textquotedblleft{}a quarrel\textquotedblright{}.

\begin{quote}
\pa{ਪੈਸੇ} \pa{ਬਾਰੇ} \pa{ਇੱਕ} \pa{ਵੱਡਾ} \pa{ਝਗੜਾ} \pa{ਹੋਇਆ।}
\end{quote}

\begin{quote}
\emph{paise bāre ikk vaḍḍā jhagṛā hoiā.}
\end{quote}

\begin{quote}
There was a big quarrel about money.
\end{quote}

A masculine noun.

\begin{cousinweb}
\textbf{\pa{ਝਗੜਾ}} (\emph{jhagṛā}) is inherited from a reconstructed \emph{*jhagaḍ}, \textquotedblleft{}quarrel\textquotedblright{}; CDIAL lists Punjabi \emph{jhagṛā}, \textquotedblleft{}quarrel\textquotedblright{}, under it (CDIAL 5321).
\end{cousinweb}

\begin{grammarlens}[title={What you've built}]
One more word for saying what happened.
\end{grammarlens}

\subsection*{Your turn}
\begin{itemize}
\raggedright
  \item \emph{Say it:} \emph{jhagṛā}
  \item \emph{Say it:} \emph{jhagṛā}, once more
  \item \emph{Say it:} \emph{paise bāre ikk vaḍḍā jhagṛā hoiā.} — There was a big quarrel about money.
  \item \emph{Recall:} say the Punjabi for a life, then the Punjabi for to fight, then say \emph{jhagṛā} again
  \item \emph{Write it:} copy \textbf{\pa{ਝਗੜਾ}} once, checking the whole word
\end{itemize}

\begin{tcolorbox}[breakable,colback=teal!4,colframe=teal!35!black,title={Before you move on}]
What does \pa{ਝਗੜਾ} mean? (\textquotedblleft{}A quarrel\textquotedblright{}.) Say it once more.
\end{tcolorbox}
\section[\texorpdfstring{\pa{ਮਾਰਨਾ} (mārnā)}{ਮਾਰਨਾ (mārnā)}]{\pa{ਮਾਰਨਾ} (mārnā) — to hit}

\label{lesson:PA-C316-maarna}



\begin{quote}
Before the new one: say the Punjabi for to fight, then the Punjabi for a quarrel.
\end{quote}

\subsection*{You'll want to know: \pa{ਮਾਰਨਾ}}
\textbf{\pa{ਮਾਰਨਾ}} — \emph{mārnā} — \textquotedblleft{}to hit\textquotedblright{}.

\begin{quote}
\pa{ਝਗੜੇ} \pa{ਵਿੱਚ} \pa{ਮੈਂ} \pa{ਮੁੰਡੇ} \pa{ਨੂੰ} \pa{ਨਹੀਂ} \pa{ਮਾਰਿਆ।}
\end{quote}

\begin{quote}
\emph{jhagṛe vich maiṁ muṇḍe nū̃ nahīṁ māriā.}
\end{quote}

\begin{quote}
In the fight I did not hit the boy.
\end{quote}

Past (I …): \pa{ਮੈਂ} \pa{ਮੱਖੀ} \pa{ਮਾਰੀ}, \emph{maiṁ makkhī mārī}.

\begin{cousinweb}
\textbf{\pa{ਮਾਰਨਾ}} (\emph{mārnā}) is inherited from Sanskrit \emph{mārayati}, \textquotedblleft{}kills\textquotedblright{}; CDIAL lists Punjabi \emph{mārnā}, \textquotedblleft{}to kill, strike\textquotedblright{}, under it (CDIAL 10066).
\end{cousinweb}

\begin{grammarlens}[title={What you've built}]
One more word for saying what happened.
\end{grammarlens}

\subsection*{Your turn}
\begin{itemize}
\raggedright
  \item \emph{Say it:} \emph{mārnā}
  \item \emph{Say it:} \emph{mārnā}, once more
  \item \emph{Say it:} \emph{jhagṛe vich maiṁ muṇḍe nū̃ nahīṁ māriā.} — In the fight I did not hit the boy.
  \item \emph{Recall:} say the Punjabi for to fight, then the Punjabi for a quarrel, then say \emph{mārnā} again
  \item \emph{Write it:} copy \textbf{\pa{ਮਾਰਨਾ}} once, checking the whole word
\end{itemize}

\begin{tcolorbox}[breakable,colback=teal!4,colframe=teal!35!black,title={Before you move on}]
What does \pa{ਮਾਰਨਾ} mean? (\textquotedblleft{}To hit\textquotedblright{}.) Say it once more.
\end{tcolorbox}
\section[\texorpdfstring{\pa{ਗਾਲ੍ਹ} (gālh)}{ਗਾਲ੍ਹ (gālh)}]{\pa{ਗਾਲ੍ਹ} (gālh) — a swear word}

\label{lesson:PA-C316-galh}



\begin{quote}
Before the new one: say the Punjabi for a quarrel, then the Punjabi for to hit.
\end{quote}

\subsection*{You'll want to know: \pa{ਗਾਲ੍ਹ}}
\textbf{\pa{ਗਾਲ੍ਹ}} — \emph{gālh} — \textquotedblleft{}a swear word\textquotedblright{}.

\begin{quote}
\pa{ਉਸ} \pa{ਨੇ} \pa{ਮੈਨੂੰ} \pa{ਗਾਲ੍ਹ} \pa{ਕੱਢੀ।}
\end{quote}

\begin{quote}
\emph{us ne mainū̃ gālh kaḍḍhī.}
\end{quote}

\begin{quote}
He swore at me.
\end{quote}

A feminine noun.

\begin{cousinweb}
\textbf{\pa{ਗਾਲ੍ਹ}} (\emph{gālh}) is inherited from Sanskrit \emph{gāli}, \textquotedblleft{}abusive speech\textquotedblright{}; CDIAL lists Punjabi \emph{gāl} and \emph{gālī}, \textquotedblleft{}abuse\textquotedblright{}, under it (CDIAL 4145).
\end{cousinweb}

\begin{grammarlens}[title={What you've built}]
One more word for saying what happened.
\end{grammarlens}

\subsection*{Your turn}
\begin{itemize}
\raggedright
  \item \emph{Say it:} \emph{gālh}
  \item \emph{Say it:} \emph{gālh}, once more
  \item \emph{Say it:} \emph{us ne mainū̃ gālh kaḍḍhī.} — He swore at me.
  \item \emph{Recall:} say the Punjabi for a quarrel, then the Punjabi for to hit, then say \emph{gālh} again
  \item \emph{Write it:} copy \textbf{\pa{ਗਾਲ੍ਹ}} once, checking the whole word
\end{itemize}

\begin{tcolorbox}[breakable,colback=teal!4,colframe=teal!35!black,title={Before you move on}]
What does \pa{ਗਾਲ੍ਹ} mean? (\textquotedblleft{}A swear word\textquotedblright{}.) Say it once more.
\end{tcolorbox}
\section[\texorpdfstring{\pa{ਸੁਲ੍ਹਾ} (sulhā)}{ਸੁਲ੍ਹਾ (sulhā)}]{\pa{ਸੁਲ੍ਹਾ} (sulhā) — a reconciliation}

\label{lesson:PA-C316-sulha}



\begin{quote}
Before the new one: say the Punjabi for to hit, then the Punjabi for a swear word.
\end{quote}

\subsection*{You'll want to know: \pa{ਸੁਲ੍ਹਾ}}
\textbf{\pa{ਸੁਲ੍ਹਾ}} — \emph{sulhā} — \textquotedblleft{}a reconciliation\textquotedblright{}.

\begin{quote}
\pa{ਸ਼ਾਮ} \pa{ਨੂੰ} \pa{ਸਾਡੀ} \pa{ਸੁਲ੍ਹਾ} \pa{ਹੋਈ।}
\end{quote}

\begin{quote}
\emph{shām nū̃ sāḍī sulhā hoī.}
\end{quote}

\begin{quote}
In the evening we made up.
\end{quote}

A feminine noun.

\begin{cousinweb}
\textbf{\pa{ਸੁਲ੍ਹਾ}} (\emph{sulhā}) is a loan, usually traced through Persian to Arabic \emph{ṣulḥ}, \textquotedblleft{}peace, reconciliation\textquotedblright{}.
\end{cousinweb}

\begin{grammarlens}[title={What you've built}]
One more word for saying what happened.
\end{grammarlens}

\subsection*{Your turn}
\begin{itemize}
\raggedright
  \item \emph{Say it:} \emph{sulhā}
  \item \emph{Say it:} \emph{sulhā}, once more
  \item \emph{Say it:} \emph{shām nū̃ sāḍī sulhā hoī.} — In the evening we made up.
  \item \emph{Recall:} say the Punjabi for to hit, then the Punjabi for a swear word, then say \emph{sulhā} again
  \item \emph{Write it:} copy \textbf{\pa{ਸੁਲ੍ਹਾ}} once, checking the whole word
\end{itemize}

\begin{tcolorbox}[breakable,colback=teal!4,colframe=teal!35!black,title={Before you move on}]
What does \pa{ਸੁਲ੍ਹਾ} mean? (\textquotedblleft{}A reconciliation\textquotedblright{}.) Say it once more.
\end{tcolorbox}
